\subsection{Every genuine length-two reference arm: the complete directory}\label{proof:R2-C034}
\paragraph{R34.1 Exact theorem, actual directions and exhaustive structural table}

On the original finite probability space put

\fitdisplay{
Jv=(Ev)1,\quad B=-K,\quad u=B\one,\quad T=J+B,\quad A=J-B,
\quad p=1-u,\quad q=1+u,
}{}
\fitdisplay{
f_n=A^{n}\one,\quad \bar f_n=T^{n}\one,\quad h_n=T f_{n-1},\qquad
\delta_i=\|B^{i}\one\|^2,\quad d=\delta_1,\quad
\theta=13/32,\quad\kappa=19/32,\quad\tau=d-\delta_2.
}{}

The receiver in the theorem is exactly the existing outward arm
\(C=f_2=A^2 1\). There is no variable receiving length in the theorem.
For every strictly increasing positive finite residual table \(\Sigma\) of
size at least two and every integer \(L\ge1\), select the actual terminal
product in the table below. Every connector edge points from terminal to
receiver, so its actual message is \(X=T^LF\). Then

\fitdisplay{
E[f_2X^4]\le1.
}{R34.1.1}

The directions use only \(\Sigma,\allowbreak  L\). In an \(f_s\) arm every edge points
toward its leaf; in an \(h_s\) arm the first edge points toward the terminal
and all remaining edges toward the leaf; in \(\bar f_s\) every edge points toward
the terminal. These statements specify each original edge, including all
degree-two subdivisions. A split of the proof by \(d\) leaves them unchanged.

\begin{description}[style=nextline,leftmargin=1.6em]
\item[{even, \((1,\allowbreak  b)\)}] Actual terminal product, all unmentioned pairs retained: \(qf_b\);\par\noindent Proof: exact Section~\ref{proof:Q5-C058}.
\item[{even, \((2,\allowbreak  b)\)}] Actual terminal product, all unmentioned pairs retained: \(h_2f_b\);\par\noindent Proof: exact Section~\ref{proof:Q5-C077}.
\item[{even, \((3,\allowbreak  b),\allowbreak  b\ge4\)}] Actual terminal product, all unmentioned pairs retained: \(f_3\bar f_b\) for \(L=1,\allowbreak  2\);\par\noindent \(h_3f_b\) for \(L\ge3\);\par\noindent Proof: §4.
\item[{even, \(s_1<t_1<\cdots<s_j<t_j\), \(s_1\ge4\)}] Actual terminal product, all unmentioned pairs retained: \(\prod h_{s_i}f_{t_i}\);\par\noindent Proof: §5.
\item[{even, \(j\ge2,\allowbreak  s_1=2\) or \(3\)}] Actual terminal product, all unmentioned pairs retained: \(\prod h_{s_i}f_{t_i}\);\par\noindent Proof: §6.
\item[{even, \(1<b_1<a_2<b_2<\cdots<a_j<b_j\), \(j\ge2\)}] Actual terminal product, all unmentioned pairs retained: \(qf_{b_1}\prod_{i=2}^j h_{a_i}f_{b_i}\);\par\noindent Proof: §7.
\item[{odd, \(2\le\ell<s_1<t_1<\cdots<s_j<t_j\), \(s_1\ge4\)}] Actual terminal product, all unmentioned pairs retained: \(\bar f_\ell\prod h_{s_i}f_{t_i}\);\par\noindent Proof: Lemma~\ref{lem:R34:nonunit-odd}, family \(\mathcal C\).
\item[{odd, \((2,\allowbreak  3,\allowbreak  c,\allowbreak  s_2,\allowbreak  t_2,\allowbreak  \ldots)\)}] Actual terminal product, all unmentioned pairs retained: \(\bar f_2f_3h_c\prod_{i=2}^j h_{s_i}f_{t_i}\);\par\noindent Proof: Lemma~\ref{lem:R34:nonunit-odd}, family \(\mathcal P\).
\item[{odd, \((1,\allowbreak  2,\allowbreak  3)\)}] Actual terminal product, all unmentioned pairs retained: \(qf_2\bar f_3\) at \(L=1\);\par\noindent \(qf_2h_3\) at \(L\ge2\);\par\noindent Proof: Section~\ref{proof:Q5-C114};\par\noindent Section~\ref{proof:Q5-C104}.
\item[{odd, \((1,\allowbreak  2,\allowbreak  4)\)}] Actual terminal product, all unmentioned pairs retained: \(qf_2h_4\);\par\noindent Proof: Section~\ref{proof:Q5-C105}.
\item[{odd, \((1,\allowbreak  2,\allowbreak  c),\allowbreak  c\ge5\)}] Actual terminal product, all unmentioned pairs retained: \(qf_2h_c\);\par\noindent Proof: §12.
\item[{odd, \((1,\allowbreak  3,\allowbreak  c),\allowbreak  c\ge4\)}] Actual terminal product, all unmentioned pairs retained: \(qf_3h_c\);\par\noindent Proof: §10.
\item[{odd, \((1,\allowbreak  b,\allowbreak  c),\allowbreak  4\le b<c\)}] Actual terminal product, all unmentioned pairs retained: \(qf_bh_c\);\par\noindent Proof: §11.
\item[{odd, \(j\ge2,\allowbreak  s_1\ge4\)}] Actual terminal product, all unmentioned pairs retained: \(q\prod h_{s_i}f_{t_i}\);\par\noindent Proof: Lemma~\ref{lem:R34:unit-multipair}, family \(A\).
\item[{odd, \(j\ge2,\allowbreak  s_1=3\)}] Actual terminal product, all unmentioned pairs retained: \(qf_3h_{t_1}\prod_{i=2}^j h_{s_i}f_{t_i}\);\par\noindent Proof: Lemma~\ref{lem:R34:unit-multipair}, family \(B\).
\item[{odd, \(j\ge2,\allowbreak  s_1=2,\allowbreak  t_1\ge4\)}] Actual terminal product, all unmentioned pairs retained: \(qf_2h_{t_1}\prod_{i=2}^j h_{s_i}f_{t_i}\);\par\noindent Proof: Lemma~\ref{lem:R34:unit-multipair}, family \(II\).
\item[{odd, \(j\ge2,\allowbreak (s_1,\allowbreak  t_1)=(2,\allowbreak  3),\allowbreak  s_2\ge5\)}] Actual terminal product, all unmentioned pairs retained: \(qf_2h_3\prod_{i=2}^j h_{s_i}f_{t_i}\);\par\noindent Proof: Lemma~\ref{lem:R34:unit-multipair}, family \(III\).
\item[{odd, \(j\ge2,\allowbreak (s_1,\allowbreak  t_1,\allowbreak  s_2)=(2,\allowbreak  3,\allowbreak  4)\)}] Actual terminal product, all unmentioned pairs retained: \(qf_2h_{t_2}h_3f_4\prod_{i=3}^j h_{s_i}f_{t_i}\);\par\noindent Proof: Lemma~\ref{lem:R34:unit-multipair}, families \(IV,V\).
\end{description}

The last row crosses the roles of the first two pairs, as prescribed by their lengths. The table is exhaustive by parity, minimum,
second minimum and, in the last odd case, the next original integer.

Arbitrary original equal-pair factors \(E_o,\allowbreak  E_t\) satisfy

\fitdisplay{
E[E_o f_2\{T^L(E_tF)\}^4]\le1.
}{R34.1.2}

Every positive length, arbitrary original repeats and unbounded finite
numbers of pairs are retained. The complete reversal uses the same result
with \(-K\), on the same measure.

\paragraph{R34.2 Genuine f2 calculus and scalar estimates}

All strict inequalities proportional to \(d\) in the proof are used for
\(d>0\). At \(d=0\), \(B\one=0\), and successive integration
of the original leaves gives every actual tree density one.

\paragraph{R34.2.1 Finite original identities}

The supplied C014 calculus gives

\fitdisplay{
0\le d\le1/3,\quad d^2\le \delta_2\le\theta d,\quad
\delta_{i+1}\le\theta \delta_i,\quad\|B\|^2\le\theta,
\quad |BV(x)|\le\|V\|.
}{}

The original positive operators have kernel in \([0,\allowbreak  2]\). The path
messages \(f_n,\allowbreak  \bar f_n\) lie in \([0,\allowbreak  2]\) and have mean at most one. One must
not extend the mean assertion to \(h_n\): \(Eh_2=1+d\).

Write

\fitdisplay{
G_n(z)=\sum_{r=0}^{n-1}z^r,\quad G_0=0,\quad
\beta_1=\beta_2=0,\quad
\beta_n=\sum_{i=1}^{\lfloor(n-1)/2\rfloor}(-1)^{i-1}\delta_i.
}{}

The exact finite positive convolution is

\fitdisplay{
1+G_n(B)u=\bar f_n+\sum_{r=3}^n\beta_r \bar f_{n-r},\qquad
\tau\le\beta_r\le d\ (r\ge3),\quad\beta_3=\beta_4=d.
}{R34.2.1}

All sums are finite; positive geometric series bound their subcollections.
At \(z=i\sqrt x\), \(0\le x\le\theta\), write \(G_n=A_n+zC_n\).
For \(n\ge2\), \(1-x\le A_n,\allowbreak  C_n\le1\); for \(n\ge3\),
\(A_n\le1-x+x^2\); for \(n\ge4\), the same holds for \(C_n\).
The exceptional \(G_1=1\) has imaginary coordinate zero. Every finite
\(G_n,\allowbreak  n=1\) or \(n\ge3\), contracts, and
\(|G_2|<19/16<6/5\). In particular, \(G_2\) is not contractive.

\paragraph{R34.2.2 Exact receiving covariance and basic products}

Directly,

\fitdisplay{
C=f_2=1-u+Bu,\quad EC=1-d,\quad
0\le C\le p+\sqrt d\le p+\tfrac65\sqrt d,
\quad C\le2-q^2/2,
}{R34.2.2}
\fitdisplay{
\sigma(C)^2=d+\delta_2-d^2\le45d/32<(6/5)^2d.
}{}

The point cap in \(2.2\) is also a direct original-row fact. For the fixed
row \(v(y)=B(x,\allowbreak  y)\), skewness gives
\(Bu(x)=\tfrac12E_{y,\allowbreak  z}[(v(y)-v(z))B(y,\allowbreak  z)]\).
Since \(|v(y)-v(z)|\le1-v(y)v(z)\) on \([-1,\allowbreak  1]^2\),
\(Bu(x)\le(1-u(x)^2)/2\), which is exactly \(C\le2-q^2/2\).

Let \(\nu\) be the original spectral measure of \(-B^2\) at \(u\):
its mass is \(d\), support is \([0,\allowbreak  \theta]\), and
\(\int x^r\,\allowbreak  d\nu=\delta_{r+1}\). If \(H=Q(B)u\) and \(EH=-\alpha\),
then the exact identity is

\fitdisplay{
\operatorname{Cov}(C,H)
 =-\int\operatorname{Re}[(1+z)Q(z)]\,d\nu-d\alpha.
}{R34.2.3}

The final signed mean product is retained in the covariance estimates.

The decisive genuine-f2 products are

\fitdisplay{
\begin{array}{c|c}
n&\operatorname{Re}(G_2G_n)\\\hline
1&1\\
2&1-x\ge\kappa\\
3&1-2x\ge3/16\\
n\ge4&\ge(1-x)^2\ge\kappa^2=361/1024.
\end{array}
}{R34.2.4}

For even \(n\), \(A_n=C_n\ge1-x\). For odd \(n\ge5\),
\(A_n=1-xC_n\) and \(C_n\le1-x+x^2\), whence
\(A_n-xC_n\ge1-2x+2x^2-2x^3\ge(1-x)^2\).
Every product in \(2.4\) is at most one. Also

\fitdisplay{
|\operatorname{Re}(G_2z^r)|=x^{\lceil r/2\rceil}\le x
\quad(r\ge2),\qquad \operatorname{Re}(G_2z)=-x,
}{}
\fitdisplay{
\operatorname{Re}(G_2(2+z))=2-x\le2,\quad
\operatorname{Re}(G_2(3+2z+z^2))=3-3x\le3.
}{R34.2.5}

For later positive additions, the actual bare covariances are

\fitdisplay{
\begin{aligned}
\operatorname{Cov}(C,q)&=-d,\\
\operatorname{Cov}(C,\bar f_2)&=-d+\delta_2-d^2<0,\\
\operatorname{Cov}(C,\bar f_3)&=-d+2\delta_2-d^2<0,\\
\operatorname{Cov}(C,\bar f_4)&=-d+2\delta_2-\delta_3+d\delta_2\le-667d/3072,\\
\operatorname{Cov}(C,\bar f_5)&=-d+2\delta_2-2\delta_3+d^2+d^3\le-337d/4608.
\end{aligned}
}{R34.2.6}

For the last two use \(v=\delta_2/d\le\theta\), \(\delta_3/d\ge v^2\),
and monotonicity of \(-1+2v-v^2+dv\) and
\(-1+2v-2v^2+d+d^2\) on the stated rectangle. Their respective
upper corners give the displayed exact rationals.

\paragraph{R34.2.3 Taylor, signed cubic means and exact bracket compensation}

For nonnegative actual \(Y\), with its actual \(m=EY\),

\fitdisplay{
E[CY^4]=(1-d)m^4+4m^3\operatorname{Cov}(C,Y)
 +E[C(Y-m)^2\{(Y+m)^2+2m^2\}].
}{R34.2.7}

If \(W=e^{-\Gamma}Y\), the first two terms scale by \(e^{-4\Gamma}\).
A pointwise coefficient \(M\) proved on \(W\) multiplies
\(e^{-2\Gamma}\operatorname{\operatorname{Var}}Y\), not \(e^{-4\Gamma}\operatorname{\operatorname{Var}}Y\).
If instead \(M\) is proved on \(Y\), all terms have the fourth-power scale.

We will use the following comparison of scalar brackets. Suppose an old bracket
is \((1-\eta d)m^4+4vm^3\), \(1\le\eta\le2\). If \(0\le m\le M\)
and

\fitdisplay{
\operatorname{Cov}(C,Y)\le v-(\eta-1)Md/4,
}{R34.2.8}

then the new mean-plus-covariance bracket is no larger: its difference is
at most \((\eta-1)d m^3(m-M)\le0\). Usually \(M=1\).
This compensates for the mean \(1-d\). When covariance is negative,
all actual cubic means or their valid lower bounds are kept; no replacement
of \(m^3\) by an upper bound is made with the wrong sign.

\paragraph{R34.2.4 Pointwise curvature and scalar certificates}

For \(x,\allowbreak  y,\allowbreak  m\ge0\), \(x+y\le H\), \(H\ge3m\), the cubic estimate is

\fitdisplay{
x(y^2+2m^2)\le4H^3/27+2Hm^2/3+m^3/2.
}{R34.2.9}

To verify it, put \(y=H-x\), then \(x=H/3+t\). For \(t\le0\) the
remaining cubic terms are nonpositive; for \(0\le t\le H/3\), complete
the quadratic square to get at most \(3m^4/(2H)\le m^3/2\).
On \(x\ge2H/3\) the polynomial is convex and its two endpoint values
are covered by the same bound using \(H^2\ge9m^2\).
Thus the curvature constants used below are:


\begingroup\small
\setlength{\tabcolsep}{3pt}
\setlength{\LTleft}{0pt}\setlength{\LTright}{\fill}
\begin{longtable}{@{}p{0.1200\linewidth}p{0.4000\linewidth}p{0.4000\linewidth}@{}}
\caption{Every genuine length-two reference arm: the complete directory: parameter values.}\label{tab:parameters:r2-c034:2}\\
\toprule
Entry & Bound for \(m\) & Valid coefficient \\
\midrule\endfirsthead
\toprule
Entry & Bound for \(m\) & Valid coefficient \\
\midrule\endhead
\bottomrule\endfoot
\(4\) & \(1\); & \(13\). \\
\addlinespace[2pt]
\(18/5\) & \(1\); & \(10\). \\
\addlinespace[2pt]
\(17/5\) & \(33/32\); & \(9\). \\
\addlinespace[2pt]
\(18/5\) & \(21/20\); & \(41/4\). \\
\end{longtable}
\endgroup


The same statement applies to \(W\) and \(EW\). All the pointwise proofs
below use the actual \(2.2\); the stronger old \((1-d)p\) bound is not assigned to \(f_2\).

If \(C_0\ge A_0/2\) and \(A_0>B_0\ge0\), then

\fitdisplay{
e^{-A_0d}+B_0d e^{-C_0d}
 \le1-(A_0-B_0)d e^{-A_0d/2}.
}{R34.2.10}

This follows from \(2\sinh(A_0d/2)\ge A_0d\). We use the elementary
finite positive-series estimates \(e<3,\allowbreak  e^2<9,\allowbreak  e^4<81\).
Other, tighter continuous certificates with their polynomial lower bounds are stated below.

\paragraph{R34.2.5 Finite masks and positive propagation}

For an actual ordered pair put

\fitdisplay{
R_{s,t}=\sum_{r=s+1}^{t}K^{r}\one.
}{R34.2.11}

Disjoint original pairs give disjoint finite intervals and a single
finite coefficient mask; no sum of separate worst-case norms is charged
once per pair. For any finite coefficient polynomial \(P\) with coefficients
in \([0,\allowbreak  1]\), split its four residue classes. Positive and negative
subcollections in each coordinate give

\fitdisplay{
|P(i\sqrt x)|^2\le\frac{1+x}{(1-x^2)^2}
 \le\frac{32768}{16245}.
}{R34.2.12}

The actual finite pair inequalities are

\fitdisplay{
h_2+f_t\le2+R_{2,t}-\kappa d,\quad
h_3+f_t\le2+R_{3,t}-(2-\theta)d,\quad
h_s+f_t\le2+R_{s,t}-16d/5\quad(s\ge4).
}{R34.2.13}

The \(s=4\) loss may be strengthened in the indicated single-pair packets.
These are finite positive-convolution facts, proved in \hyperref[part:R25.3]{the finite-mask calculation of Section~\ref{proof:R2-C025}} and the
corresponding C020/C022 packets. AM--GM applies to their actual nonnegative
left sides. A pair product is bounded by its corresponding exponential.

For \(D_i=EB^iR\) the exact Duhamel identity is

\fitdisplay{
T^rR=B^rR+\sum_{i=0}^{r-1}D_i \bar f_{r-1-i}.
}{R34.2.14}

Every upper used below is obtained by positive propagation of a true first
row. Signed \(D_i\) remain signed unless explicitly replaced by positive
parts in a positive message coefficient. The displayed \(\gamma\)'s
record every finite convolution loss, including the one needed for the
actual mean. Precursor and output norms are different quantities.

\paragraph{R34.3 Exact genuine-f2 boundary rows}




\begingroup\small
\setlength{\tabcolsep}{3pt}
\setlength{\LTleft}{0pt}\setlength{\LTright}{\fill}
\begin{longtable}{@{}p{0.2200\linewidth}p{0.7000\linewidth}@{}}
\caption{Every genuine length-two reference arm: the complete directory: parameter values.}\label{tab:parameters:r2-c034:3}\\
\toprule
Entry & Actual object and exact usable scope \\
\midrule\endfirsthead
\toprule
Entry & Actual object and exact usable scope \\
\midrule\endhead
\bottomrule\endfoot
Section~\ref{proof:Q5-C058} & \(E[f_2\{T^L(qf_b)\}^4]\le1\), every \(b\ge2,\allowbreak  L\ge1\). \\
\addlinespace[2pt]
Section~\ref{proof:Q5-C077} & \(E[f_2\{T^L(h_2f_b)\}^4]\le1-d/10000\), every \(b\ge3,\allowbreak  L\ge1\). \\
\addlinespace[2pt]
Section~\ref{proof:Q5-C101} & \(E[f_2\bar f_n^4]\le1-d/20\), every \(n\ge3\). \\
\addlinespace[2pt]
Section~\ref{proof:Q5-C114} & \(E[f_2\{T(qf_2\bar f_3)\}^4]\le1-d/256\). \\
\addlinespace[2pt]
Section~\ref{proof:Q5-C104} & \(E[f_2\{T^L(qf_2h_3)\}^4]\le1-d/500\), every \(L\ge2\). \\
\addlinespace[2pt]
Section~\ref{proof:Q5-C105} & \(E[f_2\{T^L(qf_2h_4)\}^4]\le1\), every \(L\ge1\). \\
\end{longtable}
\endgroup


Section~\ref{proof:Q5-C101}'s bare \(n=3,\allowbreak  4,\allowbreak  5\) proofs are stronger than the uniform \(d/20\),
and its exact Section~\ref{proof:Q5-C099} input covers every \(n\ge6\). Bare \(n=1,\allowbreak  2\) are not
contractions and will be paid separately whenever they arise. The
C090/C092/Section~\ref{proof:Q5-C101} larger-terminal thresholds enter in their stated ranges.

\paragraph{R34.4 The single pair (3, b) under the genuine f2 reference}

This section independently changes the receiver in the C032 finite proof. Its low-length direction is \(F=f_3\bar f_b\); its \(L\ge3\)
direction is \(F=h_3f_b\). Here \(b\ge4\).

\paragraph{R34.4.1 Length one and the full all-k precursor}

Put \(n=b-3\), \(G_n=A_0+zC_0\), and

\fitdisplay{
R=2Bu+B^3G_n(B)u,\quad\Gamma=d+\beta_b\ge51d/32,
\quad m_0=1-\Gamma/2,\quad k=m_0^{-1}\in[1,3/2].
}{}

The exact coordinates of (R/u), interpreted as a polynomial and not a
pointwise quotient, are \(x^2C_0+zr\), \(r=2-xA_0\). The finite first row
in C032 §§4--5, derived from the actual pair sum, gives

\fitdisplay{
e^\Gamma TF\le Y_1=q+kBR-c,\qquad c=3d/5,
}{}
\fitdisplay{
ER\le-51d/32,\quad D_1=EBR\le0,\quad ER^2\le21d/16,
\quad\|u+kBR\|^2\le d-2\delta_2.
}{R34.4.1}

For every \(k\in[1,\allowbreak  3/2]\), not merely its endpoints, the output symbol

\fitdisplay{
Q_k=1-kxr+zkx^2C_0
}{}

satisfies

\fitdisplay{
\operatorname{Re}(G_2Q_k)=1-kxr-kx^3C_0
 \ge1-2kx\ge1-3x\ge1-4x.
}{R34.4.2}

The signed mean product in \(2.3\) is favorable here. Equations \(4.1\),
\(4.2\), the exact \(EC=1-d\), and \(C\le p+\sqrt d\) meet every premise of
the receiver-independent C032 length-one curvature/scalar certificate
(coefficient 10). Thus its \(d/10\) deficit is retained. In particular no
bound \(|G_2|\le1\) has been inserted.

\paragraph{R34.4.2 Length two: actual negative constant, covariance and cubic mean}

The true precursor is

\fitdisplay{
V=-c\,1+u+kBR,\quad
Y_2=q+\epsilon+BV,\quad\epsilon=-c+kD_1.
}{}

It is not ((1-c)u+kBR). The entire constant cross is

\fitdisplay{
\begin{aligned}
\|V\|^2
&=\|u+kBR\|^2+c^2-2ckD_1\\
&\le d-2\delta_2+9d^2/25+(9/5)d\,\delta_3\\
&\le d-(281/160)\delta_2+9d^2/25
 \le d-(1117/800)d^2\le d.                           
\end{aligned}
}{R34.4.3}

This uses \(D_1=-\int x^2C_0\,\allowbreak  d\nu\), \(\delta_3\le\theta \delta_2\),
\(d\le1/3\), \(\delta_2\ge d^2\); its cross term has the positive sign it
actually has. Consequently \(Y_2\le q+\epsilon+\sqrt d\le q+\sqrt d\).

Write \(a_0=1-c\). The output symbol, distinct from the precursor, is

\fitdisplay{
Q_2=(a_0-kx^3C_0)+z(1-kxr).
}{}

Its genuine-f2 product is at least

\fitdisplay{
4/5-x+2x^2-2x^3>7/12.
}{R34.4.4}

The derivative is \(-6(x-1/3)^2-1/3<0\), and its value at the larger
endpoint \(5/12\) is \(2531/4320>7/12\). The possibly adverse signed
mean product is at most \(3d^2/40\). Hence

\fitdisplay{
\operatorname{Cov}(C,Y_2)\le-7d/12+3d^2/40
 \le-67d/120\le-3d/10.
}{R34.4.5}

The actual mean lies in ([1-8d/5, 1-21d/40]), in particular above
\(7/15\). The required function is now

\fitdisplay{
g(\mu)=(1-d)\mu^4-(6/5)d\mu^3.
}{}

On this actual interval,
\(g'(\mu)/\mu^2=4(1-d)\mu-(18/5)d\ge56/45-6/5=2/45>0\).
It is therefore legitimate to use the upper mean endpoint in the
negative cubic bracket. At that endpoint its exponential rate is at least
\(33/8\); combined with \(4\Gamma\ge51d/8\) it is stronger than \(52/5\).
The full output variance is at most ([\(13/12-27d/25\)]d). The genuine
pointwise bound above gives coefficient 9. Thus the exact continuous
certificate reconstructed in C032 and A016 applies:

\fitdisplay{
e^{-52d/5}+\left(39d/4-243d^2/25\right)e^{-51d/16}
 \le1-3d/320,\qquad0\le d\le1/3.
}{R34.4.6}


\paragraph{R34.4.3 Every length at least three}

For \(d\le1/8\), the canonical \(h_3f_b\) packet is

\fitdisplay{
R=\tau u+B^3G_{b-3}(-B)u,
\quad\Gamma\ge35d/16,\quad U=\tau-x^2C_0-xzA_0,
\quad|U|<9/20.
}{}

The same finite first row has \(\sigma<55d/448\),
\(D\le47d/80\), \(\epsilon=D+D_1\in[-d^2,\allowbreak  11d/16]\), and
\(D+(D_1)_+\le11d/16\). The exact losses include
\(\gamma_3\ge d,\allowbreak  \gamma_4\ge243d/256,\allowbreak  \gamma_n\ge227d/448\)
for \(n\ge5\). The actual means and norms are

\fitdisplay{
EY_3\le1-97d/80,\quad EZ_{L\ge4}\le1-13d/16,
\quad|P_n|\le1\ (n\ge3),
}{}
\fitdisplay{
|Q_2|<7/5,\quad|Q_3|\le1+4d/5,\quad
|Q_n|\le1+16d/15\ (n\ge4).
}{}

For \(P_n=G_n+z^nU\), the genuine receiving real product is positive
for every \(n\ge3\): its minimum baseline is \(3/16\), whereas
\((6/5)(9/20)\theta^{3/2}<(6/5)(9/20)(4/15)=18/125<3/16\).
Positive \(\epsilon G_{n-2}\) is favorable; negative (epsilon) costs
at most \(d^3\), since all products \(2.4\) are between zero and one.
The \(Dz^{n-2}\) and the retained (h(1+q)) costs remain exactly those
of \(2.5\). The full signed mean terms therefore preserve the upper bounds
\(\operatorname{Cov}(C,\allowbreak  Y_3)\le9d/50+3d^2/8\) and
\(\operatorname{Cov}(C,\allowbreak  Z_{L\ge4})\le d^2\).

The actual pointwise upper \(C\le p+\sqrt d\) meets the coefficients
9 and 8. The complete low-domain scalar certificates are consequently

\fitdisplay{
e^{-13d}+9d e^{-111d/40}\le1-63d/2048,
\quad
e^{-12d}+8d e^{-269d/120}\le1-21d/640.
}{R34.4.7}

For \(d\ge1/8\), the same canonical actual pair satisfies
\(h_3f_b\le e^{-d}\). The centered finite proof uses
\(\beta_*=\beta_{b-1}\); \(b=4\) is its actual short endpoint, not an
unavailable geometric tail. Hence

\fitdisplay{
E[f_2\{T^L(h_3f_b)\}^4]
 \le e^{-4d}E[f_2\bar f_L^4]\le e^{-4d}(1-d/20),\quad L\ge3.
}{}

This completes (3, b), with a conservative common \(d/1000\) deficit.

\paragraph{R34.5 Every even table of minimum at least four}

Let \(F=\prod_{i=1}^j h_{s_i}f_{t_i}\), \(j\ge1,\allowbreak  s_1\ge4\).
The entire finite C022 mask has

\fitdisplay{
F\le e^{R-\Gamma},\quad\Gamma\ge16jd/5,
\quad R=BV_0,\quad\|V_0\|^2<d/7,
\quad ER^2<d/18,\quad |R|<1/4.
}{}

For the single pair starting at four, the stronger loss is (69d/20).
The actual moments satisfy

\fitdisplay{
D_0\in[-d/5,d/12],\quad D_1\in[-d/12,d/5],\quad
D_i\le d/12\ (i\ge2).
}{}

Since \(e^R\le1+R+5R^2/8\), take \(D=11d/72\ge D_0+5ER^2/4\).
Positive propagation gives

\fitdisplay{
e^\Gamma T^LF\le Y_L=\bar f_L+B^LR+D \bar f_{L-1}
                +\sum_{i=1}^{L-1}D_i \bar f_{L-1-i}.
}{R34.5.1}

The finite loss \(\gamma_r=\beta_r-D_{r-1}\ge d/2\) for every \(r\ge3\).
The \(r=1,\allowbreak  2\) terms are paid separately in long rows. The successive
spectral tail constants for \(BR,\allowbreak  B^2R,\allowbreak  \ldots,\allowbreak  B^6R\), in units of
\(\sqrt d\), are respectively

\fitdisplay{
5/32,\quad1/10,\quad1/16,\quad1/25,\quad1/40,\quad1/60.
}{R34.5.2}

They follow directly from \(\|B\|^2\le\theta\) and the same finite
mask, so cover every subsequent finite power as well.

\paragraph{R34.5.1 Low domain \(0<d\le1/16\), with all new receiving fees}

At \(L=1\), the actual mean satisfies
\(1+5d/72\le m\le1+127d/360\le21/20\), and
\(\operatorname{\operatorname{Var}}Y\le1369d/1024\). By \(2.3\) or Cauchy on the tail,

\fitdisplay{
\operatorname{Cov}(C,Y)\le-d+(6/5)(5/32)d=-13d/16.
}{}

The negative cubic therefore contributes at most (-13d/4), since
\(m\ge1\). The actual pointwise sum is

\fitdisplay{
C+Y+m\le3+(31/25)\sqrt d+(91/180)d<18/5,
}{}

so coefficient \(41/4\) is legal. The mean exponent is \(223/18>12\),
and the combined variance/covariance coefficient is

\fitdisplay{
(41/4)(1369/1024)-13/4=42817/4096<11.
}{}

Thus \(2.10\) applies with \((A_0,\allowbreak  B_0,\allowbreak  C_0)=(12,\allowbreak  11,\allowbreak  64/5)\).

At \(L=2\),

\fitdisplay{
m\le1-203d/360,\quad m\ge1-313d/288\ge9/10,
\quad\operatorname{\operatorname{Var}}Y\le169d/100,
}{}
\fitdisplay{
\operatorname{Cov}(C,Y)\le-\kappa d+(6/5)(1/10)d
=-379d/800\le-9d/20.
}{}

Hence \(4m^3\operatorname{Cov}(C,\allowbreak  Y)\le-6561d/5000\le-5d/4\),
with the lower, not upper, mean used. The genuine pointwise sum is

\fitdisplay{
C+Y+m\le3+(69/32)\sqrt d-7d/120<18/5.
}{}

The coefficient is 10, and the scalar triple is

\fitdisplay{
(A_0,B_0,C_0)=(289/18,313/20,64/5).
}{}

Its gap is \(73/180\); since \(A_0d/2<1\), the deficit exceeds
\(73d/540>d/100\).

For \(L=3,\allowbreak  4,\allowbreak  5\), use the actual positive upper from the full finite
convolution, with all negative coefficients treated as in \(5.1\). It has

\fitdisplay{
m\le1-7d/5,\quad\operatorname{\operatorname{Var}}Y\le121d/100,
\quad\operatorname{Cov}(C,Y)\le(6/5)(11/10)d=33d/25.
}{}

The pointwise sums are bounded by
\(3+(37/16)\sqrt d-(609/360)d\) at \(L=3\), and by
\(3+(21/10)\sqrt d\) at \(L=4,\allowbreak  5\), all below \(18/5\).
For \(L\ge6\), pay the two exceptional finite coefficients by
\(E_0=127d/180\), retain the loss \(h=d/2\), and use the same shifted
finite polynomial as C022 §7. Then

\fitdisplay{
m\le1-8d/5,\quad\operatorname{\operatorname{Var}}Y\le121d/100,
\quad C+Y+m\le3+(81/40)\sqrt d+(37/180)d+d^2/2<18/5.
}{}

The same Cauchy covariance bound applies. All \(L\ge3\) therefore have
coefficient \(10\) and common scalar triple
\((19,\allowbreak  18,\allowbreak  64/5)\); the actual remainder coefficient is only
\(10(121/100)+4(33/25)=869/50<18\). Every low row has deficit at least
\(d/100\).

\paragraph{R34.5.2 High domain, including the two bare short obstructions}

For \(d\ge1/16\), the same actual entire pair product satisfies
\(F\le e^{-2d}\), by the finite centered C022 pair bound. For \(L\ge3\)
use the exact bare genuine-f2 theorem. At \(L=1\), the direct coefficient
11 with \(C\le p+3/4\), \(Eq=1\), \(\operatorname{Cov}(C,\allowbreak  q)=-d\), gives

\fitdisplay{
E[Cq^4]\le1+6d,
\quad e^{-8d}E[Cq^4]\le1-2d e^{-8d}<1-2d/27.
}{}

At \(L=2,\allowbreak  d\in[1/16,\allowbreak  1/8]\), coefficient 10 is valid because
\(C\le p+9/20\); \(2.6\) gives

\fitdisplay{
E[C\bar f_2^4]\le(1-d)^5+10d(45/32-d)-(19/8)(1-d)^3d
 \le1+(30603/4096)d<1+8d.
}{R34.5.3}

For the second inequality use \((1-d)^5\le1-5d+10d^2\) and
\((1-d)^3\ge(7/8)^3\); the \(d^2\) terms cancel. Comparison with
\(e^{8d}\ge1+8d\) leaves more than \(d/100\) after scaling.

At \(L=2,\allowbreak  d\ge1/8\), coefficient 12 is valid under \(C\le p+3/4\).
Discarding the negative covariance,

\fitdisplay{
e^{-8d}E[C\bar f_2^4]
 \le e^{-8d}(1-d)^5+12d(45/32-d)e^{-8d}
 \le239349/262144<14/15\le1-d/5.
}{R34.5.4}

Indeed the first term is at most \(50421/262144\), using \(e^{-1}<3/8\).
The derivative of \(d(45/32-d)e^{-8d}\) is negative on ([\(1/8,\allowbreak  1/3\)]):
the derivative of its bracket (\(45/32-2d-8d(45/32-d))\) is negative there,
and its value at \(1/8\) is negative. Its second term is at most
\(1476/2048\). Thus every minimum-at-least-four even table has deficit
at least \(d/100\), including all single pairs and all positive \(L\).

For clarity, the two bare \(\bar f_2\) curvature constants above have complete
one-variable certificates. The actual bound
\(\bar f_2+Eb_2\le13/4-p\) follows from
\(\sqrt d-d\le1/4\). On the lower strip the polynomial is
\fitdisplay{
\Phi(p)=(p+9/20)[(13/4-p)^2+225/128]\le10.
}{}
On \(0\le p\le3/2\), \(\Phi''\le-3\),
\(\Phi(1)=25317/2560<99/10\), and
\(|\Phi'(1)|=189/640<3/10\); its quadratic-tangent excess is below
\(3/200\). Beyond \(3/2\), its derivative is negative and decreasing.
On the upper strip use
\fitdisplay{
\Psi(p)=(p+3/4)[(13/4-p)^2+49/32]\le12.
}{}
For \(p\le1\), \(\Psi''\le-5\),
\(\Psi(3/4)=747/64,\allowbreak  \Psi'(3/4)=9/32\); its excess is at most
\(81/10240\). On \([1,\allowbreak  2]\), the convex quadratic derivative has two
negative endpoint values, so it remains negative. These are continuous
proofs of the exact coefficients used in \(5.3\)--\(5.4\).

\paragraph{R34.6 Even multipairs of minimum two or three}

Let \(j\ge2\), \(s_1\in\{2,3\}\), and retain the actual
canonical terminal \(F=\prod_{i=1}^j h_{s_i}f_{t_i}\), with its
strict original arm order. Lemma~\ref{lem:R25:even-mask}, at the
genuine receiving length \(a=2\), gives
\[
 \E[f_2(T^LF)^4]\le1-d/100
 \qquad\text{for every integer }L\ge1.
\]
The shared proof retains the full finite mask, signed first constants,
the constant vector in the length-two precursor, and all actual
long mean indices. Its genuine \(G_2\) receiving branch is proved
in paragraph~\ref{part:R25:genuine-two}. The original directions
and the decoration bound are therefore exactly those in the
reference directory.

\paragraph{R34.7 All minimum-one even multipair tables}

The fixed product is

\fitdisplay{
F=qf_{b_1}\prod_{i=2}^j h_{a_i}f_{b_i},\quad
1<b_1<a_2<b_2<\cdots<a_j<b_j,\quad j\ge2.
}{R34.7.1}

For classes I, II and IIIa, Section~\ref{proof:R2-C028} supplies the
finite masks, positive first rows, norms and scalar certificates;
the genuine \(f_2\) receiving products are checked below.
For class IIIb, Lemma~\ref{lem:R28:global-mask} applies with the actual
receiver length \(a=2\), retaining the stronger deficit \(2d/405\).

Every higher-index \(D_i\) bound in this section is ONE-SIDED. Negative
\(D_i\) values are retained in the finite formulas.

Throughout this section, \(D_i=EB^iR\) belongs to its one actual finite
mask, and positive propagation is exactly
\fitdisplay{
Y_L=\bar f_L+B^LR+D \bar f_{L-1}+\sum_{i=1}^{L-1}D_i \bar f_{L-1-i}.
}{}
The symbols \(P_n=G_n+z^nU,\allowbreak  \ Q_n=P_n+\epsilon G_{n-2}+Dz^{n-2}\),
\(\epsilon=D+D_1\), and every finite \(\gamma_n\) use the algebraic identities in
\hyperref[part:R25:finite-propagation]{Section~\ref{proof:R2-C025}},
with this same \(U,\allowbreak  R,\allowbreak  D,\allowbreak  D_i\); the mask bounds are those of this section.
Thus all displayed mean bounds retain the extra finite index rather than setting a spectral mean to zero.

\paragraph{R34.7.1 Classes I and II}

Class I is \(b_1\ge3\); class II is \(b_1=2,\allowbreak  a_2\ge4\).
On \(d\le1/64\), write the one actual finite mask as

\fitdisplay{
R=zP_0(-z)u=Uu,\quad P_0(-z)=U_0-zV_0,\quad
c_x=(1-x^2)^{-1},\quad |U|<1.
}{}

The actual mask, before coordinate bounds, is
\fitdisplay{
P_0(w)=G_{b_1-1}(w)+
 \sum_{i=2}^j\sum_{r=a_i+1}^{b_i}w^{r-2}.
}{}
Its exponent intervals are disjoint. Together with
\(R=zP_0(-z)u\), this specifies both classes. Its complete coordinate bounds are

\fitdisplay{
\begin{array}{c|cc}
&U_0&V_0\\\hline
\mathrm I &[1-xc_x,c_x]&[1-xc_x,c_x]\\
\mathrm {II}&[1-x^3c_x,1+x^2c_x]&[-xc_x,x^2c_x].
\end{array}
}{}

The losses satisfy
\(\Gamma_{\mathrm I}\ge607d/160+16(j-2)d/5\) and
\(\Gamma_{\mathrm {II}}\ge16(j-1)d/5\).
The true first-row constants are \(D=3d/8\) and \(D=0\), respectively;
\(\epsilon=D+D_1\in[-d/8,\allowbreak  3d/8]\) or ([-d/12, d/5]).
The finite moments obey \(0\le D_2\le d/2,\allowbreak  D_3\le d/5\),
\(D_i\le d/10\) for \(i\ge4\), and every retained \(\gamma\ge3d/8\).
The resulting exact mean/variance table is


\begingroup\small
\setlength{\tabcolsep}{3pt}
\setlength{\LTleft}{0pt}\setlength{\LTright}{\fill}
\begin{longtable}{@{}p{0.1200\linewidth}p{0.4000\linewidth}p{0.4000\linewidth}@{}}
\caption{Every genuine length-two reference arm: the complete directory: parameter values.}\label{tab:parameters:r2-c034:6}\\
\toprule
Entry & Mean upper (I / II where different) & Variance upper \\
\midrule\endfirsthead
\toprule
Entry & Mean upper (I / II where different) & Variance upper \\
\midrule\endhead
\bottomrule\endfoot
1 & \((1+3d/8) / (1+d/5)\); & \((27d/20)\). \\
\addlinespace[2pt]
2 & \((1-d/8) / (1-3d/10)\); & \((13d/8)\). \\
\addlinespace[2pt]
3 & \((1-37d/40)\); & \((51/40)^2d\). \\
\addlinespace[2pt]
\(\ge4\) & \((1-3d/4)\); & \((36d/25)\). \\
\end{longtable}
\endgroup


The new length-one product has
\(Q_1=(1-xU_0)+zxV_0\). Its genuine real product is at least
\(1-x/(1-x)\ge6/19\) in I, and at least
\(\kappa-\theta^3/\kappa>0\) in II. The latter's possible signed
mean cost is at most \(d^2/5\); both covariances are therefore nonpositive.
At length two both full coordinates are positive, so the old
\((1-x)Q_e-xQ_o\) lower bound improves to \(Q_e-xQ_o\), with favorable
signed mean. Its covariance remains nonpositive.

For \(P_n=G_n+z^nU\), \(n\ge4\), the product is at least
\(\kappa^2-(6/5)\theta^2>0\). For \(n=3\) it is at least
\(3/16-(6/5)(4/15)=-53/400>-1/3\).
Equation \(2.4\) pays every signed (epsilon) term, and \(2.5\) every
\(D\) term. The actual translated mean products then preserve

\fitdisplay{
\operatorname{Cov}(C,Y_3)\le d/3+3d^2/10,\qquad
\operatorname{Cov}(C,Z_{L\ge4})\le d^2.
}{}

The original coefficient 8 is valid with \(2.2\). The full continuous
low-domain scalar rows are

\fitdisplay{
e^{-64d/5}+(54/5)d e^{-32d/5},\quad
e^{-15d}+13d e^{-6d},\quad
e^{-14d}+(53/4)d e^{-6d};
}{R34.7.2}

Each is at most \(1-d/4\) on \(0\le d\le1/64\), by the
continuous payments following (R28.6.3) in
Section~\ref{proof:R2-C028}. For \(d\ge1/64\), every later actual pair is at most
\(e^{-d/2}\). Preserve that factor:

\fitdisplay{
F\le e^{-(j-1)d/2}qf_{b_1},\qquad
E[C(T^LF)^4]\le e^{-2d}\le1-6d/5.
}{R34.7.3}

The last inequality follows from \(1-e^{-2d}\ge2d/(1+2d)\ge6d/5\).
Merely discarding these pairs would prove a weaker conclusion and would
not justify the strict margin in \(7.3\).

\paragraph{R34.7.2 Class III, two pairs}

Here \(F=qf_2h_3f_b\), \(b\ge4\). On \(d\le1/8\),

\fitdisplay{
R=[\tau+z+z^3G_{b-3}(-z)]u,\quad\Gamma\ge35d/16,\quad D=0.
}{}

Writing \(G_{b-3}=A_0+zC_0\), the actual short packet is

\fitdisplay{
\epsilon=D_1\in[-d^2,\delta_3],\quad \delta_2-\delta_3\le D_2\le \delta_2,
\quad D_3\le\theta d^2,\quad D_i\le d/10\ (i\ge4),
}{}
\fitdisplay{
\gamma_3\ge\kappa d,\quad\gamma_4\ge243d/256,
\quad\gamma_n\ge7d/16\quad(n\ge5).
}{}

The two actual output symbols are

\fitdisplay{
Q_1=(1-x+x^2A_0)+z(\tau-x^2C_0),
\quad
Q_2=(1-\tau x+x^3C_0)+z(1-x+x^2A_0).
}{R34.7.4}

The first genuine product is at least \(\kappa-\theta/8\); its only
adverse mean cost is at most \(\theta^2d^2\). Thus its covariance is
nonpositive. The second has positive coordinates, so its old lower
product \(1/4\) is preserved, giving covariance at most (-d/4).
This is the sign needed in the tight second-length cubic bracket.
The actual mean is between \(55/64\) and (1-d/2), the output variance
is at most \(11d/8-\kappa^2d^2\), and the entire constant-containing
precursor is at most \(21\sqrt d/20\). These are distinct bounds.

For the later \(P_n\), use \(|U|<31/40\):

\fitdisplay{
\operatorname{Re}(G_2P_3)\ge3/16-(6/5)(31/40)(4/15)>-4/15,
}{}
\fitdisplay{
\operatorname{Re}(G_2P_n)\ge\kappa^2-(6/5)(31/40)\theta^2>0
\quad(n\ge4).
}{}

Positive (epsilon) is favorable; its negative part costs at most
\(d^3\). Full translated covariance bounds remain \(4d/15+d^3\)
at \(L=3\), and \(3d^2/5\) for longer lengths. The other actual data are

\fitdisplay{
\begin{array}{c|cc}
L&\text{mean upper}&\text{variance upper}\\\hline
1&1+D_1\le1+\delta_3\le1+\theta^2d&d+\theta d^3\\
3&1-4d/3&169d/144\\
\ge4&1-9d/8&(1+d/5)^2d.
\end{array}
}{}

The genuine \(C\) satisfies the exact pointwise premises for coefficients
8 at lengths 1, 2 and 9 later. All low-domain scalar certificates are

\fitdisplay{
e^{-9d}+8d e^{-4d},
}{}
\fitdisplay{
e^{-247d/20}+(11d-361d^2/128)e^{-35d/8},
}{}
\fitdisplay{
e^{-13d}+(169/16)d e^{-35d/8},\qquad
e^{-13d}+9d e^{-15d/4}.
}{R34.7.5}

Their respective R2-C028 continuous deficits are at least
(d/4, 37d/640, 25d/128, 35d/64). For \(d\ge1/8\), retain the same
actual inequality \(h_3f_b\le e^{-d}\) and compare with the exact
Section~\ref{proof:Q5-C058} \(qf_2\) row. This gives \(E[C(T^LF)^4]\le e^{-4d}\) for every \(L\).

\paragraph{R34.7.3 Class III, at least three pairs, on the whole domain}

The actual terminal is
\[
 F=qf_2h_3f_{b_2}\prod_{i=3}^j h_{a_i}f_{b_i},
 \qquad 4\le b_2<a_3<b_3<\cdots<a_j<b_j,\quad j\ge3.
\]
Its finite polynomial is
\[
 P_0(w)=1+w^2G_{b_2-3}(w)
       +\sum_{i=3}^j\sum_{r=a_i+1}^{b_i}w^{r-2}.
\]
The original intervals are disjoint, so every coefficient is zero or
one, with \((P_0)_0=(P_0)_2=1\) and \((P_0)_1=0\).
The actual pair inequalities give
\[
 F\le e^{R-\Gamma},\qquad R=BP_0(-B)u+\tau u,\qquad
 \Gamma=3d-2\delta_2+\frac{16}{5}(j-2)d
       \ge\frac{431}{80}d.
\]
Lemma~\ref{lem:R28:global-mask}, in its genuine \(a=2\) case, yields
\(\E[f_2(T^LF)^4]\le1-2d/405\) for every integer \(L\ge1\),
throughout \(0\le d\le1/3\). Together with the preceding classes,
this proves the minimum-one even multipair family with deficit
\(d/3000\).

\paragraph{R34.8 Nonunit odd terminals}
\label{part:R34:nonunit-odd}

\begin{lemma}\label{lem:R34:nonunit-odd}
For every integer \(L\ge1\), the following actual terminals satisfy
\(\E[f_2(T^LF)^4]\le1-d/100\):
\[
\begin{array}{c|l|l}
 &\text{original residual lengths}&F\\\hline
\mathcal C&2\le\ell<s_1<t_1<\cdots<s_j<t_j,\quad s_1\ge4
 &\bar f_\ell\prod_{i=1}^j h_{s_i}f_{t_i}\\
\mathcal P&(2,3,c,s_2,t_2,\ldots,s_j,t_j),\quad c\ge4
 &\bar f_2 f_3h_c\prod_{i=2}^j h_{s_i}f_{t_i}.
\end{array}
\]
Here \(j\ge1\), all listed lengths increase strictly, every connector
edge points toward the receiver, and the existing receiving arm has
length two and points outward.
\end{lemma}
\begin{proof}
Use \(C=f_2\), with mass \(1-d\), pointwise bound
\(C\le p+\sqrt d\), and \(\sigma(C)<(6/5)\sqrt d\).
Leaf integration handles \(d=0\). First suppose \(0<d\le1/8\).

The actual finite packets. For \(\mathcal C\), put
\(\mu_\ell=\E\bar f_\ell\in[0,1]\).
The finite anchored calculation \eqref{eq:R20:rows} gives
\(F\le\bar f_\ell e^{R-\Gamma}\) and the positive upper
\begin{equation}\label{eq:R34:nonunit-propagation}
 e^\Gamma T^LF\le Y_L=\bar f_n+\mu_\ell B^LR+D\bar f_{L-1}
     +\mu_\ell\sum_{i=1}^{L-1}D_i\bar f_{L-1-i},
 \quad n=\ell+L,\quad D_i=\E B^iR.
\end{equation}
Write \(D=c_0d\). Its linked parameters are
\[
\begin{array}{c|c|c|c|c|c}
 &\Gamma/d\ge g&\|R\|/\sqrt d<\rho&c_0&(D_1)_+/d\le c_1
 &(D_i)_+/d\le c_2,\ i\ge2\\\hline
I:\ j=1,s_1=4&69/20&1/5&72/125&1/6&7/100\\
II:\ j=1,s_1\ge5&16/5&3/20&3/5&1/30&1/12\\
III:\ j\ge2&16j/5&1/4&5/6&1/5&1/12.
\end{array}
\]
In row \(I\), necessarily \(\ell=2,3\), and the actual
\(\beta_{\ell+1}=d\) is retained. In every row \(D\) is the
upper for the anchored first constant, including the actual
\(\bar f_\ell\) mixed terms; no larger anchor variance is combined
with a different row's smaller mask.
For \(n\le5\), raise coefficients only before nonnegative paths in
\eqref{eq:R34:nonunit-propagation}. The mean losses at
\(L=1,2,3\) are respectively
\[
 k=2-c_0-c_1,\qquad
 k=5/2-c_0-c_1-c_2,\qquad
 k=11/4-c_0-c_1-2c_2,
\]
and \(\E Y_L\le1-kd\).
The norm bounds are \(1+\sqrt\theta\rho\),
\(1+\theta\rho+c_0d\), and
\(1+\theta^{3/2}\rho+(6c_0/5+c_1)d\), in units of \(\sqrt d\).
Thus the standard deviation is at most \(7\sqrt d/6\) for
single pairs and \(5\sqrt d/4\) for multipairs.

For \(n\ge6\), the same finite convolution gives
\eqref{eq:R20:long}, with
\[
\begin{gathered}
 Y=1-h(4-d)+[G_n-h(3+2z+z^2)](B)u+\mu_\ell B^LR+E,\\
 (h/d,E/d)=(2/5,0),(1/2,81/80),(3/8,59/48).
\end{gathered}
\]
The exceptional original coefficient at \(\ell+1\) is paid by
\(E\), and the extra actual mean coefficient is retained. Hence
\(\E Y\le1-5h+3hd+E\), giving mean losses
\(37/20,13/10,97/192\). The finite shifted-symbol inequalities
\eqref{eq:R18:shiftfour}--\eqref{eq:R18:shifttwo} apply to
\(n\ge6\) and \(n-1\ge5\). The long standard deviation is at
most \((1+\rho\sqrt\theta)\sqrt d\). Consequently all single
rows have \(k\ge5/4\), variance at most \(49d/36\); all
multiple rows have \(k\ge97/192\), variance at most \(25d/16\).

For \(\mathcal P\), the actual packets
\eqref{eq:R13:terminal}, \eqref{eq:R14:envelope}, and
\eqref{eq:R18:old} use
\[
 R_0=r_0(B)u,\quad r_0=-\beta+x^2V+zxU,\quad
 \beta=\beta_{c-1},\quad U\in[1-x,1],\quad V\in[0,1].
\]
The endpoint \(c=4\) is \((U,V)=(1,0)\).
The signed real coordinate gives the useful common improvement
\[
 |r_0|^2=(-\beta+x^2V)^2+x^3U^2
 \le d^2+x^4+x^3\le1/64+\theta^4+\theta^3<1/9.
\]
Thus \(\|R_0\|<\sqrt d/3\). For \(j\ge2\), add the one
complete later mask \(S\), with \(\|S\|<3\sqrt d/20\), so
\(\|R_0+S\|<(29/60)\sqrt d\).
The first constants remain \(D=3d/5\) and \(3d/2\), and the
losses remain \(g=35/16\) and \(431/80\), respectively.
Their finite propagation is \eqref{eq:R34:nonunit-propagation}
with \(n=L+2\) and coefficient one before \(B^LR,D_i\).
The actual \(\bar f_2=1+u+Bu\) is retained in this first row.
The same finite identities give the following mean/variance data:
\[
\begin{array}{c|cc|cc}
 &\multicolumn{2}{c|}{j=1}&\multicolumn{2}{c}{j\ge2}\\
 L&k&\operatorname{Var}Y/d&k&\operatorname{Var}Y/d\\\hline
1&51/40&1&41/120&121/100\\
2&17/10&1&2/3&2209/1600\\
3&2&81/64&1&25/16\\
\ge4&37/20&1849/1600&7/24&121/100.
\end{array}
\]
For \(L\ge4\), use the distinct positive bounds
\eqref{eq:R18:longsingle}--\eqref{eq:R18:longmulti}, with
\(h=2d/5,d/3\). The multirow's negative coefficient
\(\gamma_3=d-D=-d/2\) is paid by \(\bar f_{L-1}\le2\),
giving the actual extra constant \(d\), rather than deleting it.

Genuine covariances. For \(\mathcal C\) with \(n\le5\), the
bare and every positively added short path have nonpositive covariance
by \(2.6\). Only the full tail is charged, at most
\((6/5)\rho\theta^{L/2}d\); this is at most \(4d/25\) for
single pairs and \(d/5\) for multipairs.
In the long rows the core product is at least \(\kappa^2-3h\),
its signed mean parameter \(\beta_{n+1}-2hd\) is positive, and
\[
 \kappa^2-3(h/d)/8>(6/5)\rho\sqrt\theta
\]
in each of the three linked rows. Thus all long covariances are
nonpositive, even when \(L=1\) and the anchor is long.

For \(\mathcal P\) at \(L=1\), the joint core
\(G_3+zr_0\) has genuine product at least
\(1-2x-x^2-x^3>-1/20\) and positive mean parameter.
Its covariance is at most \(d/20\); adding \(BS\) costs at
most \(3d/25\). At \(L=2\), the single core
\(G_4+z^2r_0-2d/5\) has positive imaginary coordinate and
genuine real product at least \((1-x)^2-2d/5-x^4\).
Its gap exceeds the whole later fee \(3d/40\). The added
\((9/10)du\) has negative covariance, so both single and multiple
covariances are nonpositive.
At \(L=3\), keep the signed constants
\[
 P=G_5-a_0+\epsilon z+z^Lr,\qquad
 (a_0,\epsilon)=
 \begin{cases}
 (7d/5-D_1,-2d/5),&j=1,\\
 (d/2-D_1,d/2),&j\ge2.
 \end{cases}
\]
The retained cores \(G_5-a_0+\epsilon z\) have mean parameters
\(\tau-2d^2/5\) and \(\tau+d^2/2\), both positive; the
complete \(z^Lr\) tails are paid separately in the following gaps:
\[
 \kappa^2-(7/5+\theta^2)/8-18/125>0,\qquad
 \kappa^2-(49/60+\theta/2)/8-24/125>0.
\]
Here \(D_1\) keeps its sign; its range gives the displayed upper
bound for \(a_0\), not an unjustified favorable \(D_1q\) term.
For \(L\ge4\), the genuine products of the retained four-path
or two-path core exceed their whole tail fees \(9d/100,3d/25\),
respectively, and their actual signed mean parameters are positive.
Thus every \(\mathcal P\) covariance at \(L\ge2\) is nonpositive.

One curvature coefficient. Every positive upper above satisfies
\(C+Y+\E Y<4\) and \(\E Y\le1\). For \(\mathcal C,n\le5\),
the bound is \(3+(5/2)\sqrt d<4\). In the long rows the
pointwise precursor norm is at most \((1+\rho)\sqrt d\), by
the second shifted-symbol inequality, giving the common upper
\(3+(9/4)\sqrt d+19d/48<4\).
For \(\mathcal P,L=1\), the improved \(R_0\) norm gives sums
at most
\[
 3+(38/15)\sqrt d,\qquad
 3+(161/60)\sqrt d+19d/120
\]
for single and multiple rows; both are below four using
\(\sqrt{1/8}<17/48\). At \(L=2\), the actual upper is at
most \(q+\sqrt d\) or \(q+83/192\); their sums are below four.
The finite pointwise forms \eqref{eq:R18:three},
\eqref{eq:R18:longsingle}, and \eqref{eq:R18:longmulti} give,
for \(L\ge3\), the common upper
\(3+(12/5)\sqrt d+3d/8<4\).
By \(2.9\), the one unscaled coefficient is therefore
\(K=64/5>683/54\).

One payment. If \(\E Y\le1-kd\),
\(\operatorname{Var}Y\le vd\), and
\(\operatorname{Cov}(C,Y)\le bd\) with \(b\ge0\), the exact
fourth identity gives
\[
 \E[C(T^LF)^4]\le e^{-Ad}+Bd\,e^{-4gd},\qquad
 A\le4g+1+4k,\quad B\ge(64/5)v+4b.
\]
The receiving mass \(1-d\) and every actual negative mean loss
are retained. The whole nonunit odd family is paid by five rows:
\[
\begin{array}{c|c|c|c|c|c|c|c}
 &L&g&k&v&b&A&B\\\hline
\mathcal C,j=1&\ge1&16/5&5/4&49/36&4/25&94/5&91/5\\
\mathcal C,j\ge2&\ge1&32/5&97/192&25/16&1/5&28&21\\
\mathcal P,j=1&1&35/16&51/40&1&1/20&74/5&13\\
\mathcal P,j=1&\ge2&35/16&17/10&81/64&0&33/2&81/5\\
\mathcal P,j\ge2&\ge1&431/80&7/24&25/16&1/5&23&21.
\end{array}
\]
Each has \(4g\ge A/2\), \(A-B\ge3/10\), and \(Ad/2<2\).
By \(2.10\) and \(e^2<9\), the low-domain deficit exceeds
\(d/30>d/100\).

For \(d\ge1/8\), the same actual pair comparisons give
\(F\le e^{-d}\bar f_\ell\) in \(\mathcal C\), and
\(F\le e^{-(2j-1)d}\bar f_2\) in \(\mathcal P\).
The propagated bare index is at least three in both cases, so
\eqref{eq:Q101:bare-all} yields at most \(e^{-4d}(1-d/20)\).
Since \(1-e^{-4d}\ge4d/(1+4d)\ge12d/7\), the same margin
holds in the high domain. Both closed endpoints and every original
direction agree across the split.
\end{proof}

\paragraph{R34.10 Unit odd triple (1, 3, c), every c and every length}
The genuine length-two extension \eqref{eq:Q175:f2-extension}
gives, in the actual direction \(F=qf_3h_c\),
\[
 \E[f_2\{T^L(qf_3h_c)\}^4]\le1-d/800
 \qquad(c\ge4,\ L\ge1).
\]
Its proof uses the genuine \(G_2\) product and receiving mass
\(1-d\), all five connector windows, and the high-domain path
comparison. Equal-pair decorations follow by positivity.

\paragraph{R34.11 Unit odd triples (1, b, c), \(4\le b<c\)}
Apply \eqref{eq:Q178:f2-extension} with the same terminal
\(F=qf_bh_c\), the same outward receiving arm of length two,
and the same terminal-to-receiver connector of length \(L\ge1\).
The moment is at most \(1-d/3000\) on \(0<d\le1/4\), and at
most one on \(1/4\le d\le1/3\). The upstream proof retains
the three nonunit \(b\) ranges and all four unit configurations,
including the stronger actual loss for \(b\ge6\).
At \(d=0\), use leaf integration; decorations use positivity.

\paragraph{R34.12 Unit odd triples (1, 2, c), \(c\ge5\)}
\label{part:R34:unit-two-triple}
Use \eqref{eq:Q184:f2-extension} with the actual terminal
\(F=qf_2h_c\), the same receiving arm \(f_2\), and the original
terminal-to-receiver connector \(T^L\), for every \(L\ge1\).
It gives \(1-d/800\) throughout the host domain, with loss \(d/600\)
for nonunit connectors. On \(d\le1/5\), the stronger unit and
connector-dependent bounds are those of
Lemma~\ref{lem:retained-reference-payments}.
The cases \(c=3,4\) retain their separate inputs in §3, including
the \(d/50\) low-domain bound of \eqref{eq:Q105:0.3}.

\paragraph{R34.13 A common estimate for unit-leading multipairs}
\label{part:R34:unit-multipair}

\begin{lemma}\label{lem:R34:unit-multipair}
Let all arm lengths be strictly increasing in their original order,
and let \(j\ge2\). The following six terminals satisfy
\(\E[f_2(T^LF)^4]\le1-d/800\) for every integer \(L\ge1\):
\[
\begin{array}{c|l|l}
\text{family}&\text{original lengths}&F\\\hline
A&s_1\ge4&q\prod_{i=1}^j h_{s_i}f_{t_i}\\
B&s_1=3, c=t_1\ge4&qf_3h_c\prod_{i=2}^j h_{s_i}f_{t_i}\\
II&s_1=2, c=t_1\ge4&qf_2h_c\prod_{i=2}^j h_{s_i}f_{t_i}\\
III&(s_1,t_1)=(2,3),\ s_2\ge5&qf_2h_3\prod_{i=2}^j h_{s_i}f_{t_i}\\
IV&(s_1,t_1,s_2)=(2,3,4),\ j=2&qf_2h_eh_3f_4,\quad e=t_2\ge5\\
V&(s_1,t_1,s_2)=(2,3,4),\ j\ge3&qf_2h_eh_3f_4\prod_{i=3}^j h_{s_i}f_{t_i}.
\end{array}
\]
Family \(A\) satisfies the stronger bound \(1-d/100\).
\end{lemma}
\begin{proof}
Write \(C=f_2\), so \(\E C=1-d\),
\(C\le p+\sqrt d\) and \(\sigma(C)<(6/5)\sqrt d\).
At \(d=0\), integrate the original leaves. Suppose henceforth \(d>0\).

The finite packets. The positive first rows for families \(A,B\)
are \eqref{eq:R21:firstrow} and \eqref{eq:R23:firstrow}.
For families \(II\)--\(V\), use
\eqref{eq:R24:mask}--\allowbreak\eqref{eq:R24:V}.
Together these give
\[
 F\le qe^{R-\Gamma},\qquad
 T(qe^R)\le\bar f_2+BR+D,\qquad D_i=\E B^iR.
\]
These are inequalities between actual messages, before a receiving
factor is inserted. In the following table \(D=\delta d\),
\(h=\eta d\), and \((D_i)_+\le\alpha_i d\) for \(i=1,2,3\).
\[
\begin{array}{c|c|c|c|c|c|c}
 &d_0&\Gamma/d\ge g&\|R\|/\sqrt d\le\rho&\delta&\eta
 &(\alpha_1,\alpha_2,\alpha_3)\\\hline
A&1/8&32/5&1/4&11/15&1/2&(1/5,1/12,1/12)\\
B&1/64&431/80&1/2&4/3&3/8&(1/2,1/3,1/4)\\
II&1/64&351/80&2/3&11/10&1/3&(33/64,1/4,1/4)\\
III&1/8&21/5&14/25&11/10&15/32&(1/2,1/12,1/10)\\
IV&1/8&21/4&1/2&11/10&1/2&(19/40,1/16,1/12)\\
V&1/8&169/20&3/5&7/4&1/3&(2/3,1/3,1/8).
\end{array}
\]
For \(III\), \(R=B^2u+S\) and
\(\|S\|<3\sqrt d/20\), whence
\(\rho\le\theta+3/20=89/160<14/25\).
For \(IV\), retain the signed crossing coordinates
\eqref{eq:R24:crosscoordinates}: \(R=(-A_c+zB_c)(B)u\),
\(0\le A_c\le19/40\), \(-x^2\le B_c\le19d/16\), and
\(A_c+B_c\ge d/2+x/16\). Family \(V\) adds its one complete
later mask to this same crossing. Thus the table preserves every
original interval and, in particular, a possibly negative \(B_c\).

For \(L=1,2\), take
\[
 Y_1=\bar f_2+BR+D,\qquad
 Y_2=\bar f_3+B^2R+Dq+D_1.
\]
For \(L=3,4\), raise only coefficients of nonnegative paths and take
\[
\begin{aligned}
 Y_3&=\bar f_4+B^3R+D\bar f_2+\alpha_1dq+\alpha_2d,\\
 Y_4&=\bar f_5+B^4R+D\bar f_3
              +\alpha_1d\bar f_2+\alpha_2dq+\alpha_3d.
\end{aligned}
\]
For every finite \(L\ge5\), take the retained four-path upper
\eqref{eq:R24:long}, also used in \eqref{eq:R21:long}:
\[
 Y_L=1-h(4-d)+[G_{L+1}-h(3+2z+z^2)](B)u+B^LR+2D.
\]
In each case \(0\le e^\Gamma T^LF\le Y_L\).
The finite packets give \(\beta_r-D_{r-2}\ge h\) at all retained
indices \(r\ge3\), including the additional mean index. Consequently
\(\E Y_L\le1-5h+3hd+2D\) for \(L\ge5\).

Genuine receiving covariances. Apply
\eqref{eq:reference-common:covariance} with \(a=2,t=0\).
For a joint core \(Q(B)u\) with mean \(-\zeta\), this is
\[
 \operatorname{Cov}(C,Q(B)u)
 =-\int\operatorname{Re}(G_2Q)\,d\nu-d\zeta.
\]
For \(B,II\), write respectively
\(r_0=-\beta+x^2V+zxU\) and
\(r_0=-\beta-xU+zxV\), with \(U\in[1-x,1]\), \(V\in[0,1]\).
The latter interval includes the actual \(B\)-endpoint \(c=4\), where \((U,V)=(1,0)\).
For \(III\), use \(R=B^2u+S\); for \(IV\), use the signed
coordinates above. At \(L=1,2\), the joint cores
\(G_{L+1}+z^Lr_0\) have positive imaginary coordinates, and hence
\(\zeta>0\). Their real products with \(G_2\) are at least
\[
\begin{array}{c|c|c|c}
 &L=1&L=2&\text{later-mask covariance fees}\\\hline
B&1-x-x^2-x^3&1-2x-x^4&(3d/25,3d/40)\\
II&1-x-x^3&1-2x&(3d/25,3d/40)\\
III&1-x+x^2&(1-x)^2&(3d/25,3d/40)\\
IV&1-x-19dx/16&1-2x-x^4&(0,0).
\end{array}
\]
Every displayed gap exceeds its fee throughout its stated domain.
The term \(Dq\) is favorable. These eight covariances are therefore
nonpositive. For \(A,V\) at \(L=1\), the bare gap \(\kappa d\)
exceeds \((6/5)\rho\sqrt\theta\,d\); for \(A,V\) at \(L=2\),
only the tail fee \((6/5)\rho\theta d\) need be retained.

At \(L=3,4\), all added path covariances are nonpositive by
\(2.6\), so only \(B^LR\) is charged. Uniformly its fee is
\((6/5)\rho\theta^{L/2}d\). Use
\(\theta^{3/2}<13/50\) and \(\theta^2<1/6\).
For \(L\ge5\), the core has real product at least
\(\kappa^2-3h\); its actual mean parameter is
\(\beta_{L+2}-2hd>0\). The full tail costs at most
\((6/5)\rho\theta^{5/2}d\), with \(\theta^{5/2}<8/75\).
The six parameter rows give
\(\kappa^2-3\eta d_0>(6/5)\rho(8/75)\).
Thus the long covariances, too, are nonpositive.

One payment. Put \(\mu=\E Y_L\), and suppose that
\[
 \mu\le1+ad,\quad \operatorname{Var}Y_L\le vd,\quad
 \operatorname{Cov}(C,Y_L)\le bd,\quad b\ge0.
\]
If the scaled curvature coefficient on \(W=e^{-\Gamma}Y_L\)
is \(K\), the exact fourth identity gives
\begin{equation}\label{eq:R34:unit-payment}
 \E[C(T^LF)^4]\le e^{-Ad}+Bd\,e^{-2gd}
 \quad\text{when}\quad
 A\le4g-\max\{0,4a-1\},\quad B\ge Kv+4br,
\end{equation}
where \(r\ge\mu^3\) is needed only when \(b>0\).
Indeed \((1-d)\mu^4\le e^{(4a-1)d}\), while the nonnegative
covariance fee and the variance fee have the common upper scale
\(e^{-2gd}\). No mean above one is discarded.

Here is sufficient finite data for this single payment. The mean and
variance columns follow by substituting the displayed positive uppers
in \eqref{eq:R21:propagation}, \eqref{eq:R23:propagation}, and
\eqref{eq:R24:propagation}, using their joint short norms and finite
shifted-symbol bounds. The covariance column was proved above.
\[
\begin{array}{c|c|c|c|c|c|c|c|c}
 &L&a&v&b&K&r&A&B\\\hline
A&\ge1&0&1681/900&1/5&13&1&128/5&251/10\\
B&\ge1&5/6&169/100&4/25&9&9/8&19&16\\
II&1&197/320&169/100&0&9&1&16&61/4\\
II&\ge2&3/5&25/16&16/75&9&9/8&16&61/4\\
III&\ge1&3/5&36/25&22/125&10&5/4&77/5&153/10\\
IV&1&23/40&25/16&0&10&1&19&16\\
IV&\ge2&0&121/64&4/25&10&1&21&20\\
V&1\le L\le4&17/12&64/25&3/10&10&5/3&29&28\\
V&\ge5&2&121/100&0&10&1&26&13.
\end{array}
\]
Here are direct norm checks for the variance column. At \(L=1,2\),
the joint cores for \(B,II,III,IV\) have norm bounds respectively
\((6/5,1),(6/5,1),(43/40,1),(5/4,17/16)\); add the single
\(B^LS\) tail and, at \(L=2\), \(Du\).
For \(A,V\), add the full tail to the bare \(G_2,G_3\) cores.
The ordinary triangle bounds on \(\sigma(Y_L)/\sqrt d\) at
\(L=3,4\) are respectively
\[
 1+(6\delta/5+\alpha_1)d+\rho\theta^{3/2},\qquad
 1+(-2+\delta+6\alpha_1/5+\alpha_2)d+\rho\theta^2.
\]
These use the bare centered norms \(6/5,1,1,1-2d\) for indices
\(2,3,4,5\). Only \(III\) needs its tighter joint cores:
at \(L=3\) the coordinates are
\((1-x+3d/5,1-x+x^2+11d/10)\), of norm at most \(1+3d/5\);
the square difference divided by \(x\) is at most
\(-1+\theta+d+121d^2/100<0\).
At \(L=4\), its positive real coordinate is at most the real
coordinate of \(G_5\), and its imaginary coordinate exceeds that
of \(G_5\) by \(3d/5\); the norm is therefore at most
\(1+(3/5)\sqrt\theta d<1+2d/5\).
Add \(\|B^3S\|<\sqrt d/25\),
\(\|B^4S\|<\sqrt d/40\) in these two rows.
Finally the long norm is at most \(1+\rho\theta^{5/2}\), by the
retained finite shifted-symbol inequality. Substitution of the six
parameter rows proves every listed variance bound; precursor and
output norms have not been identified.
For the four positive-covariance groups with \(a>0\), the displayed
\(r\)'s follow from \((1+ad_0)^3\); all other charged groups have
\(\mu\le1\). Every listed \(A,B\) satisfies
\eqref{eq:R34:unit-payment}.

Curvature. In \(A,B,II\), direct path-row bounds give respectively
\[
 C+Y_L+\mu\le
 3+\tfrac{12}{5}\sqrt d+\tfrac23d<4,\quad
 3+\tfrac{38}{15}\sqrt d+\tfrac{13}{3}d<\tfrac{17}{5},\quad
 3+\tfrac83\sqrt d+\tfrac{1189}{320}d<\tfrac{17}{5}.
\]
Also \(\mu\le1,33/32,33/32\), respectively. Thus \(2.9\)
gives \(K=13,9,9\), on \(Y_L\) and hence on its contraction \(W\).
For \(III,IV\), the positive rows satisfy
\[
 Y_L+\mu\le(1+\delta d)q+1+\tfrac85\sqrt d+\tfrac94d,
 \quad0\le\delta\le11/10,\quad\mu\le1+3d/5.
\]
For \(V\), they satisfy
\(Y_L+\mu\le(1+3d)q+1+2\sqrt d+21d/4\),
\(\mu\le1+2d\). These bounds follow from the same positive
coefficients displayed above, including the long excess \(2D\).
In both groups \(\E W\le1\), and the coefficient of \(q\) in
\(C+W+\E W\), using \(C\le2-q+\sqrt d\), is nonpositive.
Put \(t=\sqrt d<3/8\), and use \(e^{-gd}\le(1+gt^2)^{-1}\).
The numerator of the gap to \(18/5\) is at least
\[
\begin{array}{c|l}
III,IV&\frac35-\frac{13}{5}t+\frac{447}{100}t^2-\frac{21}{5}t^3
 \ \ge\ \frac35-\frac{13}{5}t+\frac{579}{200}t^2\\
V&\frac35-3t+\frac{827}{100}t^2-\frac{169}{20}t^3
 \ \ge\ \frac35-3t+\frac{4081}{800}t^2.
\end{array}
\]
The last two quadratics have negative discriminants
\(-47/250,-3243/1000\). Therefore \(C+W+\E W<18/5\),
and \(2.9\) gives \(K=10\) with the variance scale \(e^{-2\Gamma}\).

For every payment row, \(A-B\ge1/10\), \(2g\ge A/2\), and
\(Ad_0/2<2\). Applying \(2.10\) and \(e^2<9\) therefore gives
the common low-domain bound \(1-d/100\).

The high domains use the same directed terminals. For \(A\),
\(d\ge1/8\) gives \(F\le e^{-4d}q\). At \(L\ge2\), use
\eqref{eq:Q101:bare-all}, whose index is \(L+1\ge3\). At \(L=1\),
\[
 \E[C\bar f_2^4]\le8\E[C\bar f_2]\le8(1-d)^2\le49/8,
 \qquad e^{-16d}\E[C\bar f_2^4]\le147/152<1-d/100,
\]
using \(e^{-2}<3/19\) and the nonpositive bare covariance.
For \(B,II\) at \(d\ge1/64\), every later pair is at most
\(e^{-8d/5}\). Use \eqref{eq:Q175:f2-extension} for \(B\), and
Section~\ref{proof:Q5-C105} or paragraph~\ref{part:R34:unit-two-triple}
for \(II\), in their same actual directions.
For \(III\) at \(d\ge1/8\), regroup the same product as
\[
 F=(qf_2h_{s_2})(h_3f_{t_2})\prod_{i=3}^j h_{s_i}f_{t_i}.
\]
Here \(s_2\ge5\), \(h_3f_{t_2}\le e^{-d}\), and the first factor
is the genuine base of paragraph~\ref{part:R34:unit-two-triple}.
For \(IV,V\), the fixed crossing
is already \((qf_2h_e)(h_3f_4)\) times the later pairs;
\(h_3f_4\le e^{-d}\). The two pair inequalities are the actual
pointwise comparisons at the end of Section~\ref{proof:R2-C024}.
These last three families have bound at most \(e^{-4d}\), and
all later factors are at most one. Thus the high domains retain the
claimed margins, including both closed splitting endpoints.
\end{proof}

\paragraph{R34.19 Supplementary finite certificates used in §§11--12}

\paragraph{R34.19.1 The stronger actual unit loss when \(b\ge 6\)}
Paragraph~\ref{part:Q178.5.4} proves, from the actual finite
alternating moments and \(\delta_3\ge\delta_2^2/d\),
\[
 \Gamma\ge4\tau+\delta_3\ge(4-4v+v^2)d
 \ge(2601/1024)d>5d/2,\qquad v=\delta_2/d.
\]
This is the stronger unit-row loss used in
\eqref{eq:Q178:f2-extension}.

\paragraph{R34.19.3 The weak-receiver coefficient nine}
Lemma~\ref{lem:reference-coefficient-nine}, case \(w\), applies whenever
\(C_0,Z\ge0\) and
\[
 C_0\le p+(6/5)\sqrt d,\quad
 Z\le(1-5d/8)q-d/2+(13/12)\sqrt d,\quad
 \E Z\le1-3d/2.
\]
It gives \(C_0[(Z+\E Z)^2+2(\E Z)^2]\le9\) throughout
\(0\le d\le1/3\). The actual \(f_2\) satisfies this pointwise
receiver condition by (R34.2.2).

\paragraph{R34.19.5 The high-domain path mean}
The stronger finite estimate \eqref{eq:compact-strong-path-means},
with \(83/32>5/2\), gives
\fitdisplay{
 \mu_n=\E f_n=\E\bar f_n\le1-5d/2\quad(n\ge4).
}{R34.19.2}
Its proof covers every finite short and long index on \(0\le d\le1/3\).
Consequently \(h_c\le2\mu_{c-1}\le2-5d\) for \(c\ge5\),
as used in the same-direction high comparisons above. At \(d=0\)
the bound follows directly from leaf integration.
